\section{Compiler Pipeline}\label{sec:compiler-pipeline}

The Motivo compiler is organized as a strictly sequential multi-pass pipeline.
Each stage transforms the representation of the musical program, moving from abstract textual intent to a
concrete binary artifact.
By decoupling these stages, the compiler catches errors early, keeps musical structure fully resolved before code
generation, and allows each module to be tested in isolation.

\subsection{Orchestration and Control Flow}\label{detail-subsec:orchestration}

The public entry point \texttt{motivo::compile} in \texttt{compiler/src/motivo/compiler.cpp} orchestrates the full
pipeline.
A single \texttt{DiagnosticsEngine} instance is created at the start and passed by reference into every stage that can
report structured errors; the engine accumulates \texttt{Diagnostic} records tagged with a \texttt{DiagnosticStage}
(\texttt{Parsing}, \texttt{Semantic}, \texttt{Lowering}, or \texttt{Output}) and an optional source location.
After each stage, \texttt{compiler.cpp} checks for stage-specific failures and returns early with the collected
diagnostics rather than proceeding on invalid input.

The orchestration sequence matches the implementation exactly:

\begin{enumerate}
  \item \texttt{parsing::parse\_stream} reads the input \texttt{FILE*} and returns a \texttt{ParseResult} wrapping an
    \texttt{ast::Program}, or reports parse errors and yields an empty result.
  \item \texttt{semantic::analyze} accepts the parsed program and produces an \texttt{AnalysisResult} containing the
    same \texttt{ast::Program} plus symbol-table and annotation side tables.
  \item \texttt{lowering::lower} consumes the \texttt{AnalysisResult} and returns a \texttt{LowerResult} holding an
    optional \texttt{ir::Program}.
  \item \texttt{midi::write\_midi} serializes the \texttt{ir::Program} to the output path; I/O failures are caught and
    reported as \texttt{DiagnosticStage::Output} errors.
\end{enumerate}

Figure~\ref{fig:compiler-pipeline-full} presents this control flow at full width, naming the C++ entry functions and
the guard conditions that terminate compilation early.

\subsection{Information Flow and Artifacts}\label{detail-subsec:information-flow-and-artifacts}

Each stage consumes a typed result object from its predecessor and produces the input expected by its successor.
Figure~\ref{fig:compiler-data-artifacts} summarizes the primary data structures and their fields; the annotated AST is
not a separate tree copy; semantic analysis attaches metadata to the existing \texttt{ast::Program} through
\texttt{AnalysisResult}.

\topic{\texttt{ParseResult} (\texttt{motivo::parsing}).}
On success, owns a \texttt{std::unique\_ptr\textless ast::Program\textgreater{}} built by the Bison parser.
The program decomposes into a \texttt{Header} (tempo and time signature), a vector of \texttt{GlobalItem} nodes, and a
vector of \texttt{TrackDefinition} nodes.
The \texttt{ok()} predicate is true when the program pointer is non-null.

\topic{\texttt{AnalysisResult} (\texttt{motivo::semantic}).}
Retains a const reference to the parsed \texttt{ast::Program} and owns private \texttt{SymbolTable} and
\texttt{Annotations} instances populated by \texttt{detail::Traversal::run()}.
Callers query resolved expression types and symbol identifiers through accessor methods such as
\texttt{get\_expression\_type} and \texttt{get\_resolved\_symbol}; lowering reads these caches instead of repeating
lookups.

\topic{\texttt{ir::Program} (\texttt{motivo::lowering}).}
A flat timeline artifact: \texttt{tempo\_bpm}, \texttt{time\_sig\_numerator}, \texttt{time\_sig\_denominator}, and a
vector of \texttt{ir::Track} records.
Each track stores an \texttt{music::Instrument} and a chronologically sorted vector of \texttt{ir::NoteEvent} structures
(\texttt{midi\_note}, \texttt{start\_beat}, \texttt{duration\_beats}, \texttt{velocity}).

\topic{MIDI output (\texttt{motivo::midi}).}
\texttt{write\_midi} does not return a structured result; it writes Standard MIDI File Format~1 bytes directly to disk
or throws an exception converted to an output-stage diagnostic.

\subsection{Transformation Stages}\label{detail-subsec:transformation-stages}

The transformation process is organized into four conceptual stages aligned with the pipeline diagram above.

\topic{Text to tree (frontend).} Flex (\texttt{parsing/detail/scanner.l}) tokenizes the character stream;
Bison (\texttt{parsing/detail/parser.y}) reduces tokens to an \texttt{ast::Program}.
The AST captures logical hierarchy (patterns, tracks, voices) without temporal resolution.

\topic{Tree to annotated tree (semantic analysis).} \texttt{detail::Traversal} performs symbol resolution,
explicit type checking, and musical constraint validation in a single walk.
Collectors register pattern symbols before bodies are visited; validators enforce operand and call-site rules; visitors
in \texttt{semantic/detail/traversal/visitors/} dispatch on \texttt{std::variant} node kinds.

\topic{Tree to timeline (lowering).} \texttt{detail::LowererContext} executes the annotated program like an
interpreter: \texttt{ast\_lowerers/} emit \texttt{NoteEvent}s, \texttt{expression\_evaluators/} resolve runtime
\texttt{ir::Value}s, and \texttt{value\_flattener.cpp} expands chords and sequences into atomic events.

\topic{Timeline to binary (backend).} \texttt{midi/write\_midi.cpp} converts absolute beat positions to MIDI
ticks, splits each note into Note-On and Note-Off messages, stable-sorts events, and writes delta-time encoded track
chunks.

\begin{figure}[p]
  \centering
  \motivoscalefig{%
    \begin{tikzpicture}[
        stage/.style={motivo/stage, minimum width=2.55cm, minimum height=1.65cm, text width=2.45cm},
        guard/.style={motivo/guard, text width=2.1cm},
        arrow/.style={motivo/arrow},
        diag/.style={motivo/panel, fill=motivoRose!10, minimum width=2.2cm, minimum height=6.8cm, align=center}
      ]
      \node[stage, fill=motivoSoft] (input) {\textbf{Input}\\\texttt{FILE*}\\source name\\output path};
      \node[stage, fill=motivoBlue!18, right=0.95cm of input] (parse)
      {\textbf{Parsing}\\\texttt{parse\_stream}\\Flex + Bison};
      \node[stage, fill=motivoCyan!22, right=0.95cm of parse] (sem)
      {\textbf{Semantic}\\\texttt{analyze}\\\texttt{Traversal::run}};
      \node[stage, fill=motivoGreen!22, right=0.95cm of sem] (low)
      {\textbf{Lowering}\\\texttt{lower}\\\texttt{LowererContext}};
      \node[stage, fill=motivoAmber!26, right=0.95cm of low] (midi)
      {\textbf{MIDI}\\\texttt{write\_midi}\\SMF Format~1};
      \node[stage, fill=motivoSoft, right=0.95cm of midi] (out) {\textbf{Output}\\\texttt{.mid}\\binary artifact};

      \draw[arrow] (input) -- (parse);
      \draw[arrow] (parse) -- (sem);
      \draw[arrow] (sem) -- (low);
      \draw[arrow] (low) -- (midi);
      \draw[arrow] (midi) -- (out);

      \node[guard, below=0.55cm of parse, align=center] {\texttt{!ok()}\\$\rightarrow$ abort};
      \node[guard, below=0.55cm of sem, align=center] {semantic\\errors $\rightarrow$ abort};
      \node[guard, below=0.55cm of low, align=center] {\texttt{!ok()}\\$\rightarrow$ abort};
      \node[guard, below=0.55cm of midi, align=center] {I/O exception\\$\rightarrow$ output error};

      \node[diag, above=0.35cm of sem, yshift=0.15cm] (diagbox) {%
        \textbf{DiagnosticsEngine}\\[0.35em]
        \scriptsize\sffamily
        shared ref\\[0.2em]
        passed to\\all stages\\[0.2em]
        \texttt{take\_diagnostics()}\\at end
      };

      \draw[motivo/arrow, dashed, draw=motivoRose!55] (diagbox.south) -- ++(0,-0.25) -| (parse.north);
      \draw[motivo/arrow, dashed, draw=motivoRose!55] (diagbox.south) -- (sem.north);
      \draw[motivo/arrow, dashed, draw=motivoRose!55] (diagbox.south) -- ++(0,-0.25) -| (low.north);
      \draw[motivo/arrow, dashed, draw=motivoRose!55] (diagbox.south) -- ++(0,-0.25) -| (midi.north);
    \end{tikzpicture}%
  }
  \caption{Full compiler pipeline orchestrated by \texttt{motivo::compile} in \texttt{compiler.cpp},
    including early-exit guards after each stage and the shared \texttt{DiagnosticsEngine} reference passed to every
  reporting site.}\label{fig:compiler-pipeline-full}
\end{figure}

\begin{figure}[htbp]
  \centering
  \motivoscalefig{%
    \begin{tikzpicture}[
        box/.style={motivo/box, minimum width=3.35cm, minimum height=1.35cm, text width=3.1cm},
        field/.style={motivo/guard, text width=3.0cm, align=center},
        arrow/.style={motivo/arrow, line width=1.4pt},
        label/.style={font=\small\sffamily\bfseries, text=motivoInk}
      ]
      \node[box, fill=motivoBlue!16] (pr) {\textbf{ParseResult}\\\texttt{unique\_ptr\textless Program\textgreater}};
      \node[box, fill=motivoCyan!20, right=1.1cm of pr] (ar) {\textbf{AnalysisResult}\\\texttt{Program
      ref}\\\texttt{SymbolTable}\\\texttt{Annotations}};
      \node[box, fill=motivoGreen!20, right=1.1cm of ar] (ir) {\textbf{ir::Program}\\header
      fields\\\texttt{vector\textless Track\textgreater}};
      \node[box, fill=motivoAmber!24, right=1.1cm of ir] (mid) {\textbf{SMF bytes}\\\texttt{MThd + MTrk}\\480 PPQ};

      \draw[arrow] (pr) -- node[label, above=0.15cm] {analyze} (ar);
      \draw[arrow] (ar) -- node[label, above=0.15cm] {lower} (ir);
      \draw[arrow] (ir) -- node[label, above=0.15cm] {write\_midi} (mid);

      \node[field, below=0.75cm of pr] (f1) {\texttt{Header}\\\texttt{globals[]}\\\texttt{tracks[]}\\hierarchical AST};
      \node[field, below=0.75cm of ar] (f2) {same AST nodes\\+ expression types\\+ resolved symbols};
      \node[field, below=0.75cm of ir] (f3) {flat timeline\\\texttt{NoteEvent}s\\no control flow};
      \node[field, below=0.75cm of mid] (f4) {relative delta-times\\per track chunk\\binary encoding};

      \draw[-{Stealth[length=2mm]}, draw=motivoInk!35, dashed] (pr.south) -- (f1.north);
      \draw[-{Stealth[length=2mm]}, draw=motivoInk!35, dashed] (ar.south) -- (f2.north);
      \draw[-{Stealth[length=2mm]}, draw=motivoInk!35, dashed] (ir.south) -- (f3.north);
      \draw[-{Stealth[length=2mm]}, draw=motivoInk!35, dashed] (mid.south) -- (f4.north);
    \end{tikzpicture}%
  }
  \caption{Stage artifacts passed between compiler modules: \texttt{ParseResult} $\rightarrow$
    \texttt{AnalysisResult} $\rightarrow$ \texttt{ir::Program} $\rightarrow$ MIDI
  file.}\label{fig:compiler-data-artifacts}
\end{figure}

\subsection{Technical Foundations}\label{subsec:implementation-foundations}

The implementation balances high-level musical expressiveness with low-level binary precision.
The discussion below covers the C++ language features that support this architecture.

\subsection{Language Choice: Modern C++ (C++23)}\label{detail-subsec:language-choice:-modern-c++-(c++23)}

The compiler targets the C++23 standard.
C++ was selected for deterministic performance during recursive unrolling, a strong static type system for AST and IR
modeling, and native support for binary serialization in the MIDI backend.

\topic{C++23 Features Used in the Architecture.}

\topic{\texttt{std::variant} as a sum type.} AST and IR nodes use \texttt{std::variant} (e.g.,
\texttt{StatementKind}, \texttt{ExpressionKind}, \texttt{ir::Value}) with \texttt{std::visit} dispatch in semantic
visitors, lowering lowerers, and expression evaluators.

\topic{Value-oriented memory management.} Recursive AST nodes use \texttt{std::unique\_ptr}; side tables
(\texttt{SymbolTable}, \texttt{Annotations}) use contiguous vectors indexed by opaque ID types
(\texttt{SymbolId}, \texttt{ScopeId}).

\topic{Standard ranges.} \texttt{std::ranges::stable\_sort} orders MIDI events; \texttt{std::ranges::any\_of}
checks diagnostic severity in \texttt{CompileResult::has\_errors()}.
